% Regional Summary Table: Key Metrics by U.S. Census Region
% Data aggregated from Papers 1-11 (state-level results rolled up to regions)
% Northeast: CT, ME, MA, NH, RI, VT, NJ, NY, PA (9 states, 75 districts)
% Midwest: IL, IN, MI, OH, WI, IA, KS, MN, MO, NE, ND, SD (12 states, 107 districts)
% South: DE, DC, FL, GA, MD, NC, SC, VA, WV, AL, KY, MS, TN, AR, LA, OK, TX (17 entities, 179 districts)
% West: AZ, CO, ID, MT, NV, NM, UT, WY, AK, CA, HI, OR, WA (13 states, 74 districts)

\begin{table}[htbp]
\centering
\caption{Regional Summary: Algorithmic Redistricting Performance by U.S. Census Region}
\label{tab:regional-summary}
\small
\begin{tabular}{lcccc}
\toprule
\textbf{Metric} & \textbf{Northeast} & \textbf{Midwest} & \textbf{South} & \textbf{West} \\
\midrule
\textit{Basic Characteristics} & & & & \\
States & 9 & 12 & 17 & 13 \\
Congressional Districts & 75 & 107 & 179 & 74 \\
Mean Population Density (per km²) & 342 & 89 & 78 & 42 \\
Mean Non-White (\%) & 28 & 24 & 38 & 35 \\
\midrule
\textit{Compactness (Polsby-Popper, mean)} & & & & \\
Algorithmic Plans & 0.38 & 0.33 & 0.29 & 0.31 \\
Enacted Plans & 0.24 & 0.21 & 0.19 & 0.22 \\
Improvement (\%) & +58 & +57 & +53 & +41 \\
\midrule
\textit{VRA Compliance (MM Districts, total)} & & & & \\
Algorithmic Plans & 8 & 18 & 84 & 27 \\
Enacted Plans & 4 & 9 & 45 & 10 \\
Surplus & +4 & +9 & +39 & +17 \\
\midrule
\textit{42\% Threshold (states meeting)} & & & & \\
Algorithmic Plans & 0 / 9 & 1 / 12 & 9 / 17 & 3 / 13 \\
Threshold Range & 18-32\% & 15-38\% & 24-61\% & 22-51\% \\
\midrule
\textit{Efficiency Gap (\%, mean 2016-2020)} & & & & \\
Algorithmic Plans & $-3.5$ & $-2.9$ & $-3.2$ & $-3.3$ \\
Enacted Plans & $+4.2$ & $+6.8$ & $+5.6$ & $+5.1$ \\
Difference (pp) & 7.7 & 9.7 & 8.8 & 8.4 \\
\midrule
\textit{Temporal Stability (2010→2020, \% tracts retained)} & & & & \\
Recursive Bisection & 82 & 79 & 78 & 81 \\
N-way Partitioning & 68 & 65 & 64 & 67 \\
Advantage (pp) & +14 & +14 & +14 & +14 \\
\bottomrule
\end{tabular}
\vspace{0.5em}

\footnotesize
\textit{Notes}: Regional classifications follow U.S. Census Bureau definitions. MM (majority-minority) districts defined as $\geq 50\%$ non-white population. 42\% threshold indicates states where total minority population $\geq 42\%$ (empirical threshold for achieving near-proportional MM representation). Efficiency gap: positive values favor Republicans, negative favor Democrats. Polsby-Popper compactness ranges 0-1 (higher = more compact). Temporal stability measures census tract retention when redistricting across decades. Data aggregated from Papers 1-11; see individual papers for state-level details. South includes DC. Improvement calculated as (Algorithmic - Enacted) / Enacted × 100.

\end{table}
